\documentclass[a4paper,10pt]{article}

\usepackage[utf8]{inputenc}
\usepackage{t1enc}

\usepackage[utf8]{inputenc}
\usepackage{t1enc}
\usepackage[spanish]{babel}
\usepackage[pdftex,usenames,dvipsnames]{color}
\usepackage[pdftex]{graphicx}
\usepackage{enumerate}
\usepackage{amsmath}
\usepackage{amsfonts}
\usepackage{amssymb}
\usepackage[table]{xcolor}
\usepackage[small,bf]{caption}
\usepackage{float}
\usepackage{subfig}
\usepackage{listings}
\usepackage{bm}
\usepackage{times}
\usepackage{verbatim}
\usepackage{moreverb}
\usepackage{fancyvrb}
% \usepackage{hyperref}
\usepackage{multirow}
\usepackage{url}
\usepackage{listings}
\lstset{breaklines=true}
\lstset{numbers=left, numberstyle=\scriptsize\ttfamily, numbersep=10pt, captionpos=b} 
\lstset{basicstyle=\small\ttfamily}
\lstset{framesep=4pt}

\begin{document}
\setcounter{secnumdepth}{5}
\setcounter{tocdepth}{5}

\begin{titlepage}
        \vfill
        \thispagestyle{empty}
        \begin{center}
                \includegraphics{./images/itba_logo.png}
                \vfill
                \Huge{Criptografia y Seguridad}\\
                \vspace{1cm}
                \Huge{Esteganografía}\\
                \vspace{1cm}
                \Huge{Trabajo Pr\'actico Especial 2}\\
        \end{center}
        \vfill
        \large{
        \begin{tabular}{lcr}
                Civile, Juan Pablo && 50453\\
                Crespo, Alvaro && 50758 \\
                Susnisky, Dario && 50592\\
        \end{tabular}
}
        \vspace{2cm}
        \begin{center}
                \large{10 de jde 2013}\\
        \end{center}
\end{titlepage}
\newpage

\setcounter{page}{1}

\section{Introducción}

En el presete Trabajo Práctico se presentan los resultados y análisis generados a partir del programa \textit{stegobmp}, implementado según la especificación
brindada por la cátedra. Dicho programa brinda la posibilidad de ocultar un archivo cualquiera en un archivo \textit{bmp}, mediante un método de estenografiado específico, 
con la posibilidad de encriptarlo. De igual forma, el programa permite recuperas el archivo oculto a partir de un archivo \textit{bmp}, que haya sido estenografiado con alguno 
de los métodos provistos.\\

Se presentan los resultados obtenidos a partir de las imágenes provistas por la cátedra, junto con el análisis de los mismos.
También se incluye el análisis de algunas cuestiones interesantes, establecidas por la cátedra.

\section{Desarrollo}

\subsection{Estegoanálisis de los archivos provistos}

\subsection{Cuestiones a analizar}

\subsubsection*{ 1) Para la implementación del programa \textit{stegobmp} se pide que la ocultación comience en el
primer componente del primer pixel. ¿Sería mejor empezar en otra ubicación? ¿Por qué?}

TODO

\subsubsection*{ 2) ¿Qué ventajas podría tener ocultar siempre en una misma componente? Por ejemplo, siempre
en el bit menos significativo de la componente azul.}

TODO

\subsubsection*{ 3) Esteganografiar un mismo archivo en un .bmp con cada uno de los tres algoritmos, y comparar
los resultados obtenidos. Hacer un cuadro comparativo de los tres algoritmos estableciendo ventajas y desventajas.}

TODO

\subsubsection*{ 4) Para la implementación del programa \textit{stegobmp} se pide que la extensión del archivo se oculte
después del contenido completo del archivo. ¿por qué no conviene ponerla al comienzo, después del tamaño de archivo?}

TODO

\subsubsection*{ 5) Explicar detalladamente el procedimiento realizado para descubrir qué se había ocultado en
cada archivo y de qué modo.}

TODO

\subsubsection*{ 6) ¿Qué se encontró en cada archivo?}

TODO

\subsubsection*{ 7) Algunos mensajes ocultos tenían, a su vez, otros mensajes ocultos. Indica cuál era ese mensaje
y cómo se había ocultado.}

TODO

\subsubsection*{ 8) Uno de los archivos ocultos era una porción de un video, donde se ve ejemplificado una manera
de ocultar información ¿cuál fue el portador?}

TODO

\subsubsection*{ 9) ¿De qué se trató el método de estenografiado que no era LSB? ¿Es un método eficaz? ¿Por qué?}

TODO

\subsubsection*{ 10) ¿Qué mejoras o futuras extensiones harías al programa \textit{stegobmp}?}

TODO

\clearpage
\appendix
\section{Anexo}

\end{document}
